\typeout{IJCAI--ECAI 26 Abstract Submission Draft}

\documentclass{article}
\pdfpagewidth=8.5in
\pdfpageheight=11in

\usepackage{ijcai26}
\usepackage{times}
\usepackage{soul}
\usepackage{url}
\usepackage[hidelinks]{hyperref}
\usepackage[utf8]{inputenc}
\usepackage[small]{caption}
\usepackage{graphicx}
\usepackage{amsmath}
\usepackage{amsthm}
\usepackage{booktabs}
\usepackage{algorithm}
\usepackage{algorithmic}
\usepackage[switch]{lineno}

\linenumbers
\urlstyle{same}

\pdfinfo{
/TemplateVersion (IJCAI.2026.0)
}

\title{African-Context Adaptation for Multimodal Misinformation Detection}
\author{Anonymous Submission}

\begin{document}

\maketitle

\\begin{abstract}
Multimodal misinformation often appears as image--text pairs in which the image is authentic but the accompanying claim reframes it in misleading ways. Although benchmark datasets such as Fakeddit have supported progress in multimodal misinformation detection, it remains unclear how well such models perform on African-context data. In this work, we study adaptation for multimodal misinformation detection in African media settings. We construct an African-context dataset of local image--text examples labeled as misinformation or likely consistent, represent both Fakeddit and African samples using CLIP-based image--text consistency features, and compare three settings: in-domain evaluation on Fakeddit, direct evaluation on African data before adaptation, and adaptation using combined Fakeddit and African training data with evaluation on an unseen African holdout.

Our preliminary results show a clear adaptation benefit. A model trained only on Fakeddit performs strongly on the original benchmark but degrades substantially on African-context data, especially in misinformation recall. After adaptation with African training examples, performance on the African holdout improves markedly, with especially strong gains in misinformation recall and misinformation F1, while performance on the original Fakeddit test set remains strong. These findings suggest that benchmark-trained multimodal misinformation detectors are insufficient on their own for African contexts, and that targeted African-context adaptation can improve robustness while preserving source-domain utility. We view this work as an early step toward more inclusive multimodal misinformation detection for underrepresented media environments.
\\end{abstract}

\end{document}
